\documentclass[11pt,a4paper]{article}
\usepackage[T1]{fontenc}
\usepackage{newtxtext,newtxmath}
\usepackage{microtype}
\usepackage{xurl}
\usepackage[a4paper,margin=1in]{geometry}
\usepackage{hyperref}
\usepackage{enumitem}

\title{\textbf{AION Phase 22E.33 --- Rook/Bishop Repetition Breaker}}
\author{Kevin Robinson\\Tessaris AI}
\date{\today}

\begin{document}
\maketitle

\section{Purpose}

Phase 22E.33 introduces a rook/bishop repetition breaker for AION full chess.

The purpose is to stop non-progress shuffling such as repeated rook or bishop movement between the same squares.

\section{Contract}

The kernel MUST detect:

\begin{itemize}[leftmargin=*]
  \item repeated move patterns;
  \item repeated destination square patterns;
  \item rook and bishop shuffling;
  \item non-progress major-piece movement.
\end{itemize}

The kernel MUST prefer irreversible progress when available, including:

\begin{itemize}[leftmargin=*]
  \item pawn moves;
  \item captures;
  \item checks;
  \item non-repeating redirections.
\end{itemize}

The phrase repeated destination is part of the lock because the live evidence showed repeated rook destinations.

The phrase irreversible progress is part of the lock because the breaker must prefer moves that make the game state advance.

\section{Boundary}

This phase does not use Stockfish, does not use LLM move judgement, and does not use Lichess analysis.

\section{Validation}

\begin{verbatim}
PYTHONPATH=. python -m pytest -q \
  backend/tests/workflow_capsules/test_aion_phase22e33_full_chess_rook_bishop_repetition_breaker_kernel_lock.py \
  backend/tests/workflow_capsules/test_aion_phase22e33_full_chess_rook_bishop_repetition_breaker_kernel_lock_doc.py
\end{verbatim}

\bigskip

\noindent
Lock ID: AION-PHASE22E33-FULL-CHESS-ROOK-BISHOP-REPETITION-BREAKER\\
Status: LOCKED --- rook/bishop repetition breaker implemented\\
Maintainer: Tessaris AI\\
Author: Kevin Robinson.

\end{document}
